
\documentclass{article}
\usepackage{colortbl}
\usepackage{makecell}
\usepackage{multirow}
\usepackage{supertabular}

\begin{document}

\newcounter{utterance}

\centering \large Interaction Transcript for game `matchit\_ascii', experiment `similar\_grid\_1', episode 3 with claude{-}opus{-}4{-}8{-}azure.
\vspace{24pt}

{ \footnotesize  \setcounter{utterance}{1}
\setlength{\tabcolsep}{0pt}
\begin{supertabular}{c@{$\;$}|p{.15\linewidth}@{}p{.15\linewidth}p{.15\linewidth}p{.15\linewidth}p{.15\linewidth}p{.15\linewidth}}
    \# & $\;$A & \multicolumn{4}{c}{Game Master} & $\;\:$B\\
    \hline

    \theutterance \stepcounter{utterance}  
    & & \multicolumn{4}{p{0.6\linewidth}}{
        \cellcolor[rgb]{0.9,0.9,0.9}{
            \makecell[{{p{\linewidth}}}]{
                \texttt{\tiny{[P1$\langle$GM]}}
                \texttt{Dalyvaujate bendrame spėjimo žaidime.} \\
\texttt{Tikslas {-} išsiaiškinti, ar šis tinklelis ir kitas paveikslėlis, kurį matau tik aš, yra vienodi.} \\
\texttt{Pirmiausia apibūdinkite savo tinklelį. Tada aš pateiksiu savo aprašymą ir galėsime vienas kitam užduoti klausimų apie tinklelius, kad išsiaiškintume, ar jie vienodi. Dabar pradėkite savo trumpą tinklelio aprašymą žodžiais "APRAŠYMAS:", po kurių rašykite aprašymą. Nieko daugiau nepridėkite. Tai yra jūsų tinklelis:} \\
\texttt{X X X X X} \\
\texttt{X ▢ ▢ ▢ ▢} \\
\texttt{X X X ▢ ▢} \\
\texttt{X ▢ ▢ ▢ ▢} \\
\texttt{X ▢ ▢ ▢ ▢} \\
            }
        }
    }
    & \\ \\

    \theutterance \stepcounter{utterance}  
    & \multicolumn{4}{p{0.6\linewidth}}{
        \cellcolor[rgb]{0.9,0.9,0.9}{
            \makecell[{{p{\linewidth}}}]{
                \texttt{\tiny{[P1$\rangle$GM]}}
                \texttt{APRAŠYMAS: Tai 5x5 tinklelis. Visa viršutinė eilutė (pirma eilutė) užpildyta X. Visas kairysis stulpelis (pirmas stulpelis) užpildytas X. Be to, trečioje eilutėje pirmi trys langeliai yra X (t. y. antrasis ir trečiasis langeliai taip pat X). Visi kiti langeliai yra tušti.} \\
            }
        }
    }
    & & \\ \\

    \theutterance \stepcounter{utterance}  
    & & & \multicolumn{2}{p{0.3\linewidth}}{
        \cellcolor[rgb]{0.9,0.9,0.9}{
            \makecell[{{p{\linewidth}}}]{
                \texttt{\tiny{[GM$|$GM]}}
                \texttt{continue} \\
            }
        }
    }
    & & \\ \\

    \theutterance \stepcounter{utterance}  
    & & \multicolumn{4}{p{0.6\linewidth}}{
        \cellcolor[rgb]{0.9,0.9,0.9}{
            \makecell[{{p{\linewidth}}}]{
                \texttt{\tiny{[GM$\rangle$P2]}}
                \texttt{Dalyvaujate bendrame spėjimo žaidime.} \\
\texttt{Tikslas {-} išsiaiškinti, ar šis tinklelis ir kitas paveikslėlis, kurį matau tik aš, yra vienodi.} \\
\texttt{Pirmiausia apibūdinkite savo tinklelį. Tada aš pateiksiu savo aprašymą ir galėsime vienas kitam užduoti klausimų apie tinklelius, kad išsiaiškintume, ar jie vienodi. Dabar pradėkite savo trumpą tinklelio aprašymą žodžiais "APRAŠYMAS:", po kurių rašykite aprašymą. Nieko daugiau nepridėkite. Tai yra jūsų tinklelis:} \\
\texttt{X ▢ ▢ ▢ ▢} \\
\texttt{X ▢ ▢ ▢ ▢} \\
\texttt{X X X ▢ ▢} \\
\texttt{X ▢ ▢ ▢ ▢} \\
\texttt{X X X X X} \\
            }
        }
    }
    & \\ \\

    \theutterance \stepcounter{utterance}  
    & & & \multicolumn{4}{p{0.6\linewidth}}{
        \cellcolor[rgb]{0.9,0.9,0.9}{
            \makecell[{{p{\linewidth}}}]{
                \texttt{\tiny{[GM$\langle$P2]}}
                \texttt{APRAŠYMAS: Tai 5×5 tinklelis. Pirmasis stulpelis (kairysis kraštas) yra visiškai užpildytas X visose penkiose eilutėse. Apatinė eilutė (5{-}oji) yra visiškai užpildyta X. 3{-}ioje eilutėje užpildyti pirmi trys langeliai (X X X), likusi dalis tuščia. 1{-}oje, 2{-}oje ir 4{-}oje eilutėse užpildytas tik pirmasis langelis. Bendrai figūra primena raidę „E“.} \\
            }
        }
    }
     \\ \\

    \theutterance \stepcounter{utterance}  
    & & & \multicolumn{2}{p{0.3\linewidth}}{
        \cellcolor[rgb]{0.9,0.9,0.9}{
            \makecell[{{p{\linewidth}}}]{
                \texttt{\tiny{[GM$|$GM]}}
                \texttt{continue} \\
            }
        }
    }
    & & \\ \\

    \theutterance \stepcounter{utterance}  
    & & \multicolumn{4}{p{0.6\linewidth}}{
        \cellcolor[rgb]{0.9,0.9,0.9}{
            \makecell[{{p{\linewidth}}}]{
                \texttt{\tiny{[GM$\rangle$P2]}}
                \texttt{Tai mano APRAŠYMASAPRAŠYMAS: Tai 5x5 tinklelis. Visa viršutinė eilutė (pirma eilutė) užpildyta X. Visas kairysis stulpelis (pirmas stulpelis) užpildytas X. Be to, trečioje eilutėje pirmi trys langeliai yra X (t. y. antrasis ir trečiasis langeliai taip pat X). Visi kiti langeliai yra tušti.} \\
\texttt{Dabar užduokite klausimą, kad sužinotumėte naujus mano tinklelio aspektus, kurie gali skirtis nuo jūsų tinklelio. Pradėkite "KLAUSIMAS:" ir nieko daugiau nepridėkite.} \\
            }
        }
    }
    & \\ \\

    \theutterance \stepcounter{utterance}  
    & & & \multicolumn{4}{p{0.6\linewidth}}{
        \cellcolor[rgb]{0.9,0.9,0.9}{
            \makecell[{{p{\linewidth}}}]{
                \texttt{\tiny{[GM$\langle$P2]}}
                \texttt{KLAUSIMAS: Ar tavo apatinė eilutė (penktoji eilutė) yra visiškai užpildyta X, ar joje užpildytas tik pirmasis langelis?} \\
            }
        }
    }
     \\ \\

    \theutterance \stepcounter{utterance}  
    & & & \multicolumn{2}{p{0.3\linewidth}}{
        \cellcolor[rgb]{0.9,0.9,0.9}{
            \makecell[{{p{\linewidth}}}]{
                \texttt{\tiny{[GM$|$GM]}}
                \texttt{continue} \\
            }
        }
    }
    & & \\ \\

    \theutterance \stepcounter{utterance}  
    & & \multicolumn{4}{p{0.6\linewidth}}{
        \cellcolor[rgb]{0.9,0.9,0.9}{
            \makecell[{{p{\linewidth}}}]{
                \texttt{\tiny{[P1$\langle$GM]}}
                \texttt{Tai mano APRAŠYMASAPRAŠYMAS: Tai 5×5 tinklelis. Pirmasis stulpelis (kairysis kraštas) yra visiškai užpildytas X visose penkiose eilutėse. Apatinė eilutė (5{-}oji) yra visiškai užpildyta X. 3{-}ioje eilutėje užpildyti pirmi trys langeliai (X X X), likusi dalis tuščia. 1{-}oje, 2{-}oje ir 4{-}oje eilutėse užpildytas tik pirmasis langelis. Bendrai figūra primena raidę „E“.} \\
\texttt{KLAUSIMAS: Ar tavo apatinė eilutė (penktoji eilutė) yra visiškai užpildyta X, ar joje užpildytas tik pirmasis langelis?Pradėkite atsakymą „ATSAKYMAS:" ir nieko daugiau nepridėkite.} \\
            }
        }
    }
    & \\ \\

    \theutterance \stepcounter{utterance}  
    & \multicolumn{4}{p{0.6\linewidth}}{
        \cellcolor[rgb]{0.9,0.9,0.9}{
            \makecell[{{p{\linewidth}}}]{
                \texttt{\tiny{[P1$\rangle$GM]}}
                \texttt{ATSAKYMAS: Mano apatinėje eilutėje (penktojoje) užpildytas tik pirmasis langelis, likusieji tušti. Tavo apatinė eilutė yra visiškai užpildyta X, o mano – ne. Taip pat mano viršutinė eilutė (pirma) yra visiškai užpildyta X, o tavo – ne. Todėl mūsų tinkleliai NĖRA vienodi.} \\
            }
        }
    }
    & & \\ \\

    \theutterance \stepcounter{utterance}  
    & & & \multicolumn{2}{p{0.3\linewidth}}{
        \cellcolor[rgb]{0.9,0.9,0.9}{
            \makecell[{{p{\linewidth}}}]{
                \texttt{\tiny{[GM$|$GM]}}
                \texttt{continue} \\
            }
        }
    }
    & & \\ \\

    \theutterance \stepcounter{utterance}  
    & & \multicolumn{4}{p{0.6\linewidth}}{
        \cellcolor[rgb]{0.9,0.9,0.9}{
            \makecell[{{p{\linewidth}}}]{
                \texttt{\tiny{[P1$\langle$GM]}}
                \texttt{Dabar užduokite klausimą, kad sužinotumėte naujus mano tinklelio aspektus, kurie gali skirtis nuo jūsų tinklelio. Pradėkite "KLAUSIMAS:" ir nieko daugiau nepridėkite.} \\
            }
        }
    }
    & \\ \\

    \theutterance \stepcounter{utterance}  
    & \multicolumn{4}{p{0.6\linewidth}}{
        \cellcolor[rgb]{0.9,0.9,0.9}{
            \makecell[{{p{\linewidth}}}]{
                \texttt{\tiny{[P1$\rangle$GM]}}
                \texttt{KLAUSIMAS: Ar tavo pirmoji (viršutinė) eilutė yra visiškai užpildyta X, ar joje užpildytas tik pirmasis langelis?} \\
            }
        }
    }
    & & \\ \\

    \theutterance \stepcounter{utterance}  
    & & & \multicolumn{2}{p{0.3\linewidth}}{
        \cellcolor[rgb]{0.9,0.9,0.9}{
            \makecell[{{p{\linewidth}}}]{
                \texttt{\tiny{[GM$|$GM]}}
                \texttt{continue} \\
            }
        }
    }
    & & \\ \\

    \theutterance \stepcounter{utterance}  
    & & \multicolumn{4}{p{0.6\linewidth}}{
        \cellcolor[rgb]{0.9,0.9,0.9}{
            \makecell[{{p{\linewidth}}}]{
                \texttt{\tiny{[GM$\rangle$P2]}}
                \texttt{ATSAKYMAS: Mano apatinėje eilutėje (penktojoje) užpildytas tik pirmasis langelis, likusieji tušti. Tavo apatinė eilutė yra visiškai užpildyta X, o mano – ne. Taip pat mano viršutinė eilutė (pirma) yra visiškai užpildyta X, o tavo – ne. Todėl mūsų tinkleliai NĖRA vienodi.} \\
\texttt{KLAUSIMAS: Ar tavo pirmoji (viršutinė) eilutė yra visiškai užpildyta X, ar joje užpildytas tik pirmasis langelis?Pradėkite atsakymą „ATSAKYMAS:" ir nieko daugiau nepridėkite.} \\
            }
        }
    }
    & \\ \\

    \theutterance \stepcounter{utterance}  
    & & & \multicolumn{4}{p{0.6\linewidth}}{
        \cellcolor[rgb]{0.9,0.9,0.9}{
            \makecell[{{p{\linewidth}}}]{
                \texttt{\tiny{[GM$\langle$P2]}}
                \texttt{ATSAKYMAS: Mano viršutinėje eilutėje (pirmoje) užpildytas tik pirmasis langelis, likusieji tušti. Tavo viršutinė eilutė yra visiškai užpildyta X, o mano – ne. Todėl mūsų tinkleliai NĖRA vienodi.} \\
            }
        }
    }
     \\ \\

    \theutterance \stepcounter{utterance}  
    & & & \multicolumn{2}{p{0.3\linewidth}}{
        \cellcolor[rgb]{0.9,0.9,0.9}{
            \makecell[{{p{\linewidth}}}]{
                \texttt{\tiny{[GM$|$GM]}}
                \texttt{continue} \\
            }
        }
    }
    & & \\ \\

    \theutterance \stepcounter{utterance}  
    & & \multicolumn{4}{p{0.6\linewidth}}{
        \cellcolor[rgb]{0.9,0.9,0.9}{
            \makecell[{{p{\linewidth}}}]{
                \texttt{\tiny{[GM$\rangle$P2]}}
                \texttt{Dabar užduokite klausimą, kad sužinotumėte naujus mano tinklelio aspektus, kurie gali skirtis nuo jūsų tinklelio. Pradėkite "KLAUSIMAS:" ir nieko daugiau nepridėkite.} \\
            }
        }
    }
    & \\ \\

    \theutterance \stepcounter{utterance}  
    & & & \multicolumn{4}{p{0.6\linewidth}}{
        \cellcolor[rgb]{0.9,0.9,0.9}{
            \makecell[{{p{\linewidth}}}]{
                \texttt{\tiny{[GM$\langle$P2]}}
                \texttt{KLAUSIMAS: Ar tavo tinklelyje yra dar kokių nors užpildytų langelių, be viršutinės eilutės, kairiojo stulpelio ir trečios eilutės pirmų trijų langelių?} \\
            }
        }
    }
     \\ \\

    \theutterance \stepcounter{utterance}  
    & & & \multicolumn{2}{p{0.3\linewidth}}{
        \cellcolor[rgb]{0.9,0.9,0.9}{
            \makecell[{{p{\linewidth}}}]{
                \texttt{\tiny{[GM$|$GM]}}
                \texttt{continue} \\
            }
        }
    }
    & & \\ \\

    \theutterance \stepcounter{utterance}  
    & & \multicolumn{4}{p{0.6\linewidth}}{
        \cellcolor[rgb]{0.9,0.9,0.9}{
            \makecell[{{p{\linewidth}}}]{
                \texttt{\tiny{[P1$\langle$GM]}}
                \texttt{ATSAKYMAS: Mano viršutinėje eilutėje (pirmoje) užpildytas tik pirmasis langelis, likusieji tušti. Tavo viršutinė eilutė yra visiškai užpildyta X, o mano – ne. Todėl mūsų tinkleliai NĖRA vienodi.} \\
\texttt{KLAUSIMAS: Ar tavo tinklelyje yra dar kokių nors užpildytų langelių, be viršutinės eilutės, kairiojo stulpelio ir trečios eilutės pirmų trijų langelių?Pradėkite atsakymą „ATSAKYMAS:" ir nieko daugiau nepridėkite.} \\
            }
        }
    }
    & \\ \\

    \theutterance \stepcounter{utterance}  
    & \multicolumn{4}{p{0.6\linewidth}}{
        \cellcolor[rgb]{0.9,0.9,0.9}{
            \makecell[{{p{\linewidth}}}]{
                \texttt{\tiny{[P1$\rangle$GM]}}
                \texttt{ATSAKYMAS: Ne, mano tinklelyje nėra kitų užpildytų langelių. Užpildyta tik viršutinė eilutė, kairysis stulpelis ir trečios eilutės pirmi trys langeliai. Mūsų tinkleliai NĖRA vienodi, nes mano viršutinė eilutė pilnai užpildyta, o tavo apatinė eilutė pilnai užpildyta.} \\
            }
        }
    }
    & & \\ \\

    \theutterance \stepcounter{utterance}  
    & & & \multicolumn{2}{p{0.3\linewidth}}{
        \cellcolor[rgb]{0.9,0.9,0.9}{
            \makecell[{{p{\linewidth}}}]{
                \texttt{\tiny{[GM$|$GM]}}
                \texttt{continue} \\
            }
        }
    }
    & & \\ \\

    \theutterance \stepcounter{utterance}  
    & & \multicolumn{4}{p{0.6\linewidth}}{
        \cellcolor[rgb]{0.9,0.9,0.9}{
            \makecell[{{p{\linewidth}}}]{
                \texttt{\tiny{[P1$\langle$GM]}}
                \texttt{Dabar užduokite klausimą, kad sužinotumėte naujus mano tinklelio aspektus, kurie gali skirtis nuo jūsų tinklelio. Pradėkite "KLAUSIMAS:" ir nieko daugiau nepridėkite.} \\
            }
        }
    }
    & \\ \\

    \theutterance \stepcounter{utterance}  
    & \multicolumn{4}{p{0.6\linewidth}}{
        \cellcolor[rgb]{0.9,0.9,0.9}{
            \makecell[{{p{\linewidth}}}]{
                \texttt{\tiny{[P1$\rangle$GM]}}
                \texttt{KLAUSIMAS: Ar tavo tinklelyje yra dar kokių nors užpildytų langelių, be kairiojo stulpelio, apatinės eilutės ir trečios eilutės pirmų trijų langelių?} \\
            }
        }
    }
    & & \\ \\

    \theutterance \stepcounter{utterance}  
    & & & \multicolumn{2}{p{0.3\linewidth}}{
        \cellcolor[rgb]{0.9,0.9,0.9}{
            \makecell[{{p{\linewidth}}}]{
                \texttt{\tiny{[GM$|$GM]}}
                \texttt{continue} \\
            }
        }
    }
    & & \\ \\

    \theutterance \stepcounter{utterance}  
    & & \multicolumn{4}{p{0.6\linewidth}}{
        \cellcolor[rgb]{0.9,0.9,0.9}{
            \makecell[{{p{\linewidth}}}]{
                \texttt{\tiny{[GM$\rangle$P2]}}
                \texttt{ATSAKYMAS: Ne, mano tinklelyje nėra kitų užpildytų langelių. Užpildyta tik viršutinė eilutė, kairysis stulpelis ir trečios eilutės pirmi trys langeliai. Mūsų tinkleliai NĖRA vienodi, nes mano viršutinė eilutė pilnai užpildyta, o tavo apatinė eilutė pilnai užpildyta.} \\
\texttt{KLAUSIMAS: Ar tavo tinklelyje yra dar kokių nors užpildytų langelių, be kairiojo stulpelio, apatinės eilutės ir trečios eilutės pirmų trijų langelių?Pradėkite atsakymą „ATSAKYMAS:" ir nieko daugiau nepridėkite.} \\
            }
        }
    }
    & \\ \\

    \theutterance \stepcounter{utterance}  
    & & & \multicolumn{4}{p{0.6\linewidth}}{
        \cellcolor[rgb]{0.9,0.9,0.9}{
            \makecell[{{p{\linewidth}}}]{
                \texttt{\tiny{[GM$\langle$P2]}}
                \texttt{ATSAKYMAS: Ne, mano tinklelyje nėra kitų užpildytų langelių. Užpildyta tik kairysis stulpelis, apatinė eilutė ir trečios eilutės pirmi trys langeliai. Mūsų tinkleliai NĖRA vienodi, nes mano apatinė eilutė pilnai užpildyta, o tavo viršutinė eilutė pilnai užpildyta.} \\
            }
        }
    }
     \\ \\

    \theutterance \stepcounter{utterance}  
    & & & \multicolumn{2}{p{0.3\linewidth}}{
        \cellcolor[rgb]{0.9,0.9,0.9}{
            \makecell[{{p{\linewidth}}}]{
                \texttt{\tiny{[GM$|$GM]}}
                \texttt{continue} \\
            }
        }
    }
    & & \\ \\

    \theutterance \stepcounter{utterance}  
    & & \multicolumn{4}{p{0.6\linewidth}}{
        \cellcolor[rgb]{0.9,0.9,0.9}{
            \makecell[{{p{\linewidth}}}]{
                \texttt{\tiny{[GM$\rangle$P2]}}
                \texttt{Dabar užduokite klausimą, kad sužinotumėte naujus mano tinklelio aspektus, kurie gali skirtis nuo jūsų tinklelio. Pradėkite "KLAUSIMAS:" ir nieko daugiau nepridėkite.} \\
            }
        }
    }
    & \\ \\

    \theutterance \stepcounter{utterance}  
    & & & \multicolumn{4}{p{0.6\linewidth}}{
        \cellcolor[rgb]{0.9,0.9,0.9}{
            \makecell[{{p{\linewidth}}}]{
                \texttt{\tiny{[GM$\langle$P2]}}
                \texttt{KLAUSIMAS: Ar tavo kairiajame stulpelyje visos penkios eilutės užpildytos X, ar kuri nors iš jų tuščia?} \\
            }
        }
    }
     \\ \\

    \theutterance \stepcounter{utterance}  
    & & & \multicolumn{2}{p{0.3\linewidth}}{
        \cellcolor[rgb]{0.9,0.9,0.9}{
            \makecell[{{p{\linewidth}}}]{
                \texttt{\tiny{[GM$|$GM]}}
                \texttt{continue} \\
            }
        }
    }
    & & \\ \\

    \theutterance \stepcounter{utterance}  
    & & \multicolumn{4}{p{0.6\linewidth}}{
        \cellcolor[rgb]{0.9,0.9,0.9}{
            \makecell[{{p{\linewidth}}}]{
                \texttt{\tiny{[P1$\langle$GM]}}
                \texttt{ATSAKYMAS: Ne, mano tinklelyje nėra kitų užpildytų langelių. Užpildyta tik kairysis stulpelis, apatinė eilutė ir trečios eilutės pirmi trys langeliai. Mūsų tinkleliai NĖRA vienodi, nes mano apatinė eilutė pilnai užpildyta, o tavo viršutinė eilutė pilnai užpildyta.} \\
\texttt{KLAUSIMAS: Ar tavo kairiajame stulpelyje visos penkios eilutės užpildytos X, ar kuri nors iš jų tuščia?Pradėkite atsakymą „ATSAKYMAS:" ir nieko daugiau nepridėkite.} \\
            }
        }
    }
    & \\ \\

    \theutterance \stepcounter{utterance}  
    & \multicolumn{4}{p{0.6\linewidth}}{
        \cellcolor[rgb]{0.9,0.9,0.9}{
            \makecell[{{p{\linewidth}}}]{
                \texttt{\tiny{[P1$\rangle$GM]}}
                \texttt{ATSAKYMAS: Mano kairiajame stulpelyje visos penkios eilutės užpildytos X. Tai sutampa su tavo tinkleliu. Tačiau mūsų tinkleliai vis tiek NĖRA vienodi, nes mano viršutinė eilutė pilnai užpildyta, o tavo apatinė eilutė pilnai užpildyta.} \\
            }
        }
    }
    & & \\ \\

    \theutterance \stepcounter{utterance}  
    & & & \multicolumn{2}{p{0.3\linewidth}}{
        \cellcolor[rgb]{0.9,0.9,0.9}{
            \makecell[{{p{\linewidth}}}]{
                \texttt{\tiny{[GM$|$GM]}}
                \texttt{continue} \\
            }
        }
    }
    & & \\ \\

    \theutterance \stepcounter{utterance}  
    & & \multicolumn{4}{p{0.6\linewidth}}{
        \cellcolor[rgb]{0.9,0.9,0.9}{
            \makecell[{{p{\linewidth}}}]{
                \texttt{\tiny{[P1$\langle$GM]}}
                \texttt{Dabar užduokite klausimą, kad sužinotumėte naujus mano tinklelio aspektus, kurie gali skirtis nuo jūsų tinklelio. Pradėkite "KLAUSIMAS:" ir nieko daugiau nepridėkite.} \\
            }
        }
    }
    & \\ \\

    \theutterance \stepcounter{utterance}  
    & \multicolumn{4}{p{0.6\linewidth}}{
        \cellcolor[rgb]{0.9,0.9,0.9}{
            \makecell[{{p{\linewidth}}}]{
                \texttt{\tiny{[P1$\rangle$GM]}}
                \texttt{KLAUSIMAS: Ar tavo trečioje eilutėje užpildyti tik pirmi trys langeliai (4{-}as ir 5{-}as tušti), ar joje yra užpildyta daugiau langelių?} \\
            }
        }
    }
    & & \\ \\

    \theutterance \stepcounter{utterance}  
    & & & \multicolumn{2}{p{0.3\linewidth}}{
        \cellcolor[rgb]{0.9,0.9,0.9}{
            \makecell[{{p{\linewidth}}}]{
                \texttt{\tiny{[GM$|$GM]}}
                \texttt{continue} \\
            }
        }
    }
    & & \\ \\

    \theutterance \stepcounter{utterance}  
    & & \multicolumn{4}{p{0.6\linewidth}}{
        \cellcolor[rgb]{0.9,0.9,0.9}{
            \makecell[{{p{\linewidth}}}]{
                \texttt{\tiny{[GM$\rangle$P2]}}
                \texttt{ATSAKYMAS: Mano kairiajame stulpelyje visos penkios eilutės užpildytos X. Tai sutampa su tavo tinkleliu. Tačiau mūsų tinkleliai vis tiek NĖRA vienodi, nes mano viršutinė eilutė pilnai užpildyta, o tavo apatinė eilutė pilnai užpildyta.} \\
\texttt{KLAUSIMAS: Ar tavo trečioje eilutėje užpildyti tik pirmi trys langeliai (4{-}as ir 5{-}as tušti), ar joje yra užpildyta daugiau langelių?Pradėkite atsakymą „ATSAKYMAS:" ir nieko daugiau nepridėkite.} \\
            }
        }
    }
    & \\ \\

    \theutterance \stepcounter{utterance}  
    & & & \multicolumn{4}{p{0.6\linewidth}}{
        \cellcolor[rgb]{0.9,0.9,0.9}{
            \makecell[{{p{\linewidth}}}]{
                \texttt{\tiny{[GM$\langle$P2]}}
                \texttt{ATSAKYMAS: Mano trečioje eilutėje užpildyti tik pirmi trys langeliai (X X X), o 4{-}as ir 5{-}as langeliai yra tušti. Tai sutampa su tavo tinkleliu. Tačiau mūsų tinkleliai vis tiek NĖRA vienodi, nes mano apatinė eilutė pilnai užpildyta, o tavo viršutinė eilutė pilnai užpildyta.} \\
            }
        }
    }
     \\ \\

    \theutterance \stepcounter{utterance}  
    & & & \multicolumn{2}{p{0.3\linewidth}}{
        \cellcolor[rgb]{0.9,0.9,0.9}{
            \makecell[{{p{\linewidth}}}]{
                \texttt{\tiny{[GM$|$GM]}}
                \texttt{continue} \\
            }
        }
    }
    & & \\ \\

    \theutterance \stepcounter{utterance}  
    & & \multicolumn{4}{p{0.6\linewidth}}{
        \cellcolor[rgb]{0.9,0.9,0.9}{
            \makecell[{{p{\linewidth}}}]{
                \texttt{\tiny{[GM$\rangle$P2]}}
                \texttt{Dabar priimkite sprendimą. Kaip manote: ar jūsų tinklelis ir kitas aprašytas tinklelis yra tokie patys? Parašykite "SPRENDIMAS: tas pats tinklelis", jei manote, kad tai tas pats tinklelis, arba "SPRENDIMAS: skirtingi tinkleliai", jei manote, kad tai skirtingi tinkleliai. Nieko daugiau nepridėkite.} \\
            }
        }
    }
    & \\ \\

    \theutterance \stepcounter{utterance}  
    & & & \multicolumn{4}{p{0.6\linewidth}}{
        \cellcolor[rgb]{0.9,0.9,0.9}{
            \makecell[{{p{\linewidth}}}]{
                \texttt{\tiny{[GM$\langle$P2]}}
                \texttt{SPRENDIMAS: skirtingi tinkleliai} \\
            }
        }
    }
     \\ \\

    \theutterance \stepcounter{utterance}  
    & & & \multicolumn{2}{p{0.3\linewidth}}{
        \cellcolor[rgb]{0.9,0.9,0.9}{
            \makecell[{{p{\linewidth}}}]{
                \texttt{\tiny{[GM$|$GM]}}
                \texttt{continue} \\
            }
        }
    }
    & & \\ \\

    \theutterance \stepcounter{utterance}  
    & & & \multicolumn{2}{p{0.3\linewidth}}{
        \cellcolor[rgb]{0.9,0.9,0.9}{
            \makecell[{{p{\linewidth}}}]{
                \texttt{\tiny{[GM$|$GM]}}
                \texttt{success} \\
            }
        }
    }
    & & \\ \\

    \theutterance \stepcounter{utterance}  
    & & \multicolumn{4}{p{0.6\linewidth}}{
        \cellcolor[rgb]{0.9,0.9,0.9}{
            \makecell[{{p{\linewidth}}}]{
                \texttt{\tiny{[P1$\langle$GM]}}
                \texttt{ATSAKYMAS: Mano trečioje eilutėje užpildyti tik pirmi trys langeliai (X X X), o 4{-}as ir 5{-}as langeliai yra tušti. Tai sutampa su tavo tinkleliu. Tačiau mūsų tinkleliai vis tiek NĖRA vienodi, nes mano apatinė eilutė pilnai užpildyta, o tavo viršutinė eilutė pilnai užpildyta.} \\
\texttt{Dabar priimkite sprendimą. Kaip manote: ar jūsų tinklelis ir kitas aprašytas tinklelis yra tokie patys? Parašykite "SPRENDIMAS: tas pats tinklelis", jei manote, kad tai tas pats tinklelis, arba "SPRENDIMAS: skirtingi tinkleliai", jei manote, kad tai skirtingi tinkleliai. Nieko daugiau nepridėkite.} \\
            }
        }
    }
    & \\ \\

    \theutterance \stepcounter{utterance}  
    & \multicolumn{4}{p{0.6\linewidth}}{
        \cellcolor[rgb]{0.9,0.9,0.9}{
            \makecell[{{p{\linewidth}}}]{
                \texttt{\tiny{[P1$\rangle$GM]}}
                \texttt{SPRENDIMAS: skirtingi tinkleliai} \\
            }
        }
    }
    & & \\ \\

    \theutterance \stepcounter{utterance}  
    & & & \multicolumn{2}{p{0.3\linewidth}}{
        \cellcolor[rgb]{0.9,0.9,0.9}{
            \makecell[{{p{\linewidth}}}]{
                \texttt{\tiny{[GM$|$GM]}}
                \texttt{continue} \\
            }
        }
    }
    & & \\ \\

    \theutterance \stepcounter{utterance}  
    & & & \multicolumn{2}{p{0.3\linewidth}}{
        \cellcolor[rgb]{0.9,0.9,0.9}{
            \makecell[{{p{\linewidth}}}]{
                \texttt{\tiny{[GM$|$GM]}}
                \texttt{success} \\
            }
        }
    }
    & & \\ \\

\end{supertabular}
}

\end{document}
